\section{Заключение}

В результате выполнения учебно-исследовательской работы были изучены методы физического и логического кодирования, которые используются в телекоммуникационных системах для передачи данных. Были проведены эксперименты по кодированию исходного сообщения с использованием различных методов, а также выполнен анализ их достоинств и недостатков.

На основании проведённого анализа \textbf{наилучшим способом кодирования} для передачи исходного сообщения является чистое физическое кодирование методами \textit{Manchester (M2)} или \textit{Differential Manchester (DM)}:
\begin{itemize}
    \item Методы M2 и DM отнесены к категории лучших, эквивалентны между собой и оптимальны по соотношению «качество передачи — стоимость реализации».
    \item Сигнал часто меняется (переход в середине каждого битового интервала), что обеспечивает надёжную синхронизацию.
\end{itemize}

Использование методов логического кодирования в рамках данной работы является \textit{нецелесообразным}:
\begin{itemize}
    \item Блочное кодирование 4B/5B увеличивает длину сообщения и требует расширения полосы пропускания.
    \item Скремблирование по стандарту V.27 ter для методов M2 и DM может лишь дать потенциально снижение средней частоты сигнала, но добавит небольшую сложность реализации.
\end{itemize}